\section{Gradient-Scale Effects in Classification}

The output-gradient weighting in EF has an important classification-specific symptom. Consider one data sample. For cross-entropy with softmax probabilities $p_n$ and one-hot label $y_n$, the output-loss gradient with respect to logits is
\[
r_n=p_n-y_n.
\]
For a confidently correct prediction with class $c$ and small $\varepsilon$,
\[
p_{n,c}=1-\varepsilon,
\]
we have
\[
\|r_n\|_2^2=\|p_n-y_n\|_2^2 \in \Theta(\varepsilon^2).
\]
Thus the EF contribution of this sample becomes small as the model becomes more confident and correct. This is coherent with the way we derived the EF: the scalar first-order loss change becomes small. In continual learning, this is undesirable: The samples that the old model classifies best may receive very low importance. EWC-DR identifies this failure mode in classification and proposes a classification-specific repair through a logit-reversal operation, literally flipping the signs of the model's logits \citep{EWC-DR}. In the present framework, their observation can be read as a symptom of the weighting issue identified in Section 8: $F_{\mathrm{EF}}$ importance depends on output-gradient scale. This evidently does not occur for $M_{\mathrm{IEF}}$, and we make this point precise in the next section. Hence, the normalization analysis yields a formal problem identification and its generalizable solution, the Improved Empirical Fisher matrix.
